%% Separately authored annex additions after exact finite-matrix proof restoration.
% Necessity: the exact source gives the generic complete tableau proof but not
% the K3, LP, and FDE specialization/decision arguments or K3--LP duality proof.
% Definition: this editorial derivative descends from historical annex lines
% 39--116 and now includes additional finite-chain, Bochvar, and Post proofs;
% it is not a literal slice.
% Standing: it is editorial annex material, not inherited source, a replacement
% for the reader-facing proof, or a silent status promotion.
% Nonredundancy: historical annex item matrix.1 compresses the source proof and
% is retired; these seven items instead prove distinct claims of retained cards.
% Use: annex/p08.tex inputs this file once; route, slice, consumer, and annex
% back-reference checks fail closed if the derivative or its dependencies drift.
\noindent\textit{Annex note.} The generic finite-matrix theorem and complete
proof remain reader-facing in exact
Section~\ref{sec:many-valued-paraconsistent-02}. The following items are
separate specializations, decision arguments, and a K3--LP duality proof.

\begin{annexitem}{annex:p8.lukasiewicz.1}{thm:p08-finite-lukasiewicz-tableau}{Finite-chain tableau adequacy and decision}
\begin{proof}
Fix \(n\geq2\).  The set \(\text{\L}_n\) is finite, and its negation and
implication operations are total.  The generic labelled-tableau theorem in the
exact source applies to every finite matrix: for each nondesignated root label
it expands a compound into exactly the tuples of immediate values producing
that label.  Induction on formula complexity proves that a branch is realizable
exactly when its labels extend to a valuation in \(\text{\L}_n\).  Hence all
nondesignated-root branches close exactly when the formula is valid; the
premise-root version gives finite-premise consequence.  Every expansion lowers
the number of unexpanded connective occurrences, and only finitely many labels
exist, so exhaustive expansion terminates.
\end{proof}
\end{annexitem}

\begin{annexitem}{annex:p8.bochvar.1}{thm:p08-bochvar-tableau}{Finite-matrix adequacy for the Bochvar language}
\begin{proof}
The expanded Bochvar language still has a finite value set and total truth
operations: the internal connectives use the infectious weak-Kleene tables and
each \(J_T,J_F,J_N\) row is a total unary operation.  Instantiating the exact
generic finite-matrix tableau with those operations therefore gives the same
branch-realizability induction as in Annex~\ref{annex:p8.lukasiewicz.1}.
Closure of every nondesignated-root tableau is equivalent to validity in the
fixed expanded matrix.
\end{proof}
\end{annexitem}

\begin{annexitem}{annex:p8.post.1}{thm:p08-post-matrix-decision}{Decision for each fixed Post matrix}
\begin{proof}
For a finite input using \(m\) propositional variables, enumerate the
\(n^m\) valuations into \(V_n\), evaluate every formula recursively under the
displayed cyclic negation and maximum operation, and reject exactly when a
valuation designates all premises but not the conclusion.  The finite search
terminates and decides the fixed-matrix consequence relation.
\end{proof}
\end{annexitem}

% BEGIN HISTORICAL DERIVATIVE 39-116
\begin{annexitem}{annex:p8.k3.1}{entry:p08-k3}{K3 finite-tableau adequacy}
\begin{proof}
The K3 matrix has a nonempty finite value set and total truth operations.  Its
designated set and the generic root convention are exactly those in the card,
so the finite-matrix theorem in Section~\ref{sec:many-valued-paraconsistent-02} specializes to the stated K3
adequacy equivalence.
\end{proof}
\end{annexitem}

\begin{annexitem}{annex:p8.k3.2}{entry:p08-k3}{Decidability of finite-premise K3 consequence}
\begin{proof}
If the finite input contains \(m\) variables, enumerate its \(3^m\) K3
valuations, recursively compute every premise and the conclusion, and accept
exactly when no valuation designates all premises while leaving the conclusion
nondesignated.  The search terminates and decides the card's consequence
relation.
\end{proof}
\end{annexitem}

\begin{annexitem}{annex:p8.lp.1}{entry:p08-lp}{LP finite-tableau adequacy}
\begin{proof}
The LP matrix has the same finite total operations as K3 and the designated set
fixed by its card.  Substituting those data into
The finite-matrix theorem in Section~\ref{sec:many-valued-paraconsistent-02} gives precisely the stated LP
adequacy equivalence.
\end{proof}
\end{annexitem}

\begin{annexitem}{annex:p8.lp.2}{entry:p08-lp}{Decidability of finite-premise LP consequence}
\begin{proof}
For an input with \(m\) variables, enumerate all \(3^m\) LP valuations and
evaluate formulas recursively.  A countervaluation is exactly one designating
every premise and leaving the conclusion at \(0\).  The finite search decides
whether such a valuation exists.
\end{proof}
\end{annexitem}

\begin{annexitem}{annex:p8.fde.1}{entry:p08-fde}{FDE finite-tableau adequacy}
\begin{proof}
The support-pair FDE matrix has four values, total De Morgan operations, and the
designated set in the card.  The finite-matrix theorem in Section~\ref{sec:many-valued-paraconsistent-02}, with
exactly those data, yields the claimed finite-premise FDE adequacy equivalence.
\end{proof}
\end{annexitem}

\begin{annexitem}{annex:p8.fde.2}{entry:p08-fde}{Finite-premise FDE decision procedure}
\begin{proof}
For \(m\) input variables, enumerate the \(4^m\) support-pair valuations and
evaluate each formula by the displayed operations.  Reject exactly when some
valuation gives positive support to every premise but not to the conclusion.
This finite exhaustive test decides FDE consequence.
\end{proof}
\end{annexitem}

\begin{annexitem}{annex:p8.matrix.2}{thm:p08-k3-lp-facts}{K3 and LP theorem sets and duality}
\begin{proof}
Giving every variable value \(\tfrac12\) in K3 gives every formula that value,
by induction on \(\neg,\land,\lor\).  Since it is nondesignated, K3 has no
tautologies.

For strong Kleene operations, simultaneous induction shows that if a formula
has crisp value \(0\) or \(1\), every replacement of middle-valued variables
by crisp values preserves that result.  Negation exchanges the crisp values;
for conjunction, crisp \(1\) requires both conjuncts crisp \(1\), while crisp
\(0\) supplies a crisp \(0\) conjunct; disjunction is dual.  Thus an LP
valuation making \(A\) nondesignated, hence value \(0\), has a classical
completion falsifying \(A\).  Conversely, every classical valuation is an LP
valuation.  Therefore LP-valid formulas are exactly the classical tautologies.

Finally, a K3 counterexample to
\(\Gamma\entails_{\mathbf K_3}\Delta\) gives every member of \(\Gamma\) value \(1\) and
every member of \(\Delta\) value \(0\) or \(\tfrac12\).  Under negation these
become, respectively, nondesignated LP conclusions and designated LP premises,
so they are exactly a counterexample to
\(\neg\Delta\entails_{\mathsf{LP}}\neg\Gamma\).  Negating the same argument
reverses the correspondence, proving the equivalence.
\end{proof}
\end{annexitem}
% END HISTORICAL DERIVATIVE 39-116
